\documentclass[tikz,border=10pt]{standalone}
\usetikzlibrary{calc}

\begin{document}
\begin{tikzpicture}
    % Define coordinates
    \coordinate (A) at (0, 5);
    \coordinate (C) at (0, 0);
    \coordinate (E) at (0, -4);
    \coordinate (B) at (-8, 2);
    
    % Calculate D and H based on geometric constraints
    % Line BC: y = -0.25x
    % H is projection of A onto BC
    \coordinate (H) at (-1.176, 0.294); 
    % D is on BC such that AD perp DE
    \coordinate (D) at (-4.45, 1.11);

    % Draw solid lines
    \draw (A) -- (B);
    \draw (A) -- (D);
    \draw (B) -- (C); % This covers B-D-H-C
    \draw (D) -- (E);
    \draw (A) -- (E); % This covers A-C-E

    % Draw dashed line
    \draw[dashed] (A) -- (H);

    % Draw right angle at H (AH perp BC)
    \draw ($(H)!0.25cm!(A)$) -- ($(H)!0.25cm!(A) + ($(H)!0.25cm!(C)$) - (H)$) -- ($(H)!0.25cm!(C)$);

    % Draw right angle at D (AD perp DE)
    \draw ($(D)!0.25cm!(A)$) -- ($(D)!0.25cm!(A) + ($(D)!0.25cm!(E)$) - (D)$) -- ($(D)!0.25cm!(E)$);

    % Labels
    \node[above right] at (A) {$A$};
    \node[left] at (B) {$B$};
    \node[right] at (C) {$C$};
    \node[left] at (D) {$D$};
    \node[below right] at (E) {$E$};
    \node[below right] at (H) {$H$};
\end{tikzpicture}
\end{document}
